% Part: first-order-logic
% Chapter: sequent-calculus
% Section: soundness-identity

\documentclass[../../../include/open-logic-section]{subfiles}

\begin{document}

\olfileid[pt]{fol}{seq}{sid}

\olsection{Correção com \usetoken{S}{identity}}

\begin{prop}
$\Log{LK}$ com sequentes iniciais e regras para a identidade é correto.
\end{prop}

\begin{proof}
Sequentes iniciais da forma ${} \Sequent \eq[t][t]$ são válidos, pois
para toda !!{structure}~$\Struct M$, $\Sat{M}{\eq[t][t]}$. (Note que
supomos que o termo $t$ seja fechado, isto é, que não contenha variáveis,
de modo que atribuições de variáveis são irrelevantes).

Suponha que a última inferência em !!a{derivation} seja $=$. Então a premissa
é $\eq[t_1][t_2], \Gamma \Sequent \Delta, !A(t_1)$ e a conclusão
é $\eq[t_1][t_2], \Gamma \Sequent \Delta, !A(t_2)$. Considere
!!a{structure}~$\Struct M$. Precisamos mostrar que a conclusão é
válida, isto é, se $\Sat{M}{\eq[t_1][t_2]}$ e $\Sat{M}{\Gamma}$, então
ou $\Sat{M}{!C}$ para algum $!C \in \Delta$ ou $\Sat{M}{!A(t_2)}$.

Pela hipótese de indução, a premissa é válida. Isso significa que, se
$\Sat{M}{\eq[t_1][t_2]}$ e $\Sat{M}{\Gamma}$, então ou (a) para algum $!C
\in \Delta$, $\Sat{M}{!C}$, ou (b) $\Sat{M}{!A(t_1)}$. No caso (a),
terminamos. Considere o caso (b).  Seja $s$ uma atribuição de variáveis com
$s(x) = \Value{t_1}{M}$.  Por \olref[syn][ass]{prop:sentence-sat-true},
$\Sat{M}{!A(t_1)}[s]$. Como $\varAssign{s}{s}{x}$, por
\olref[syn][ext]{prop:ext-formulas}, $\Sat{M}{!A(x)}[s]$. Como
$\Sat{M}{\eq[t_1][t_2]}$, temos $\Value{t_1}{M} = \Value{t_2}{M}$,
e portanto $s(x) = \Value{t_2}{M}$.  Aplicando
\olref[syn][ext]{prop:ext-formulas} novamente, temos também
$\Sat{M}{!A(t_2)}[s]$.  Por \olref[syn][ass]{prop:sentence-sat-true},
$\Sat{M}{!A(t_2)}$.
\end{proof}

\end{document}
